\section{Giới thiệu}

% Slide 3
\begin{frame}[t]{\textsl{PAPI - Chỉ số Hiệu quả Quản trị và Hành chính công cấp tỉnh}}
  \begin{columns}[T,onlytextwidth]
    \begin{column}{0.56\textwidth}
      \small
      PAPI đo lường trải nghiệm, cảm nhận và mức độ hài lòng của người dân đối với quản trị, hành chính công và dịch vụ công tại địa phương.

      \vspace{0.25cm}
      \begin{itemize}
        \setlength{\itemsep}{6pt}
        \item[-] Khảo sát người dân, không phải dữ liệu hành chính.
        \item[-] Điểm được tổng hợp theo tỉnh và lĩnh vực.
        \item[-] Điểm cao hơn phản ánh trải nghiệm tích cực hơn trong mẫu khảo sát.
        \item[-] Đồ án không sử dụng dữ liệu từng cá nhân.
      \end{itemize}
    \end{column}
    \begin{column}{0.40\textwidth}
      \centering
      \begin{tikzpicture}[
        node distance=0.18cm,
        survey-stage/.style={rectangle, rounded corners, draw=Logo1, fill=Logo1!7,
          align=center, minimum height=0.66cm, text width=3.7cm, font=\scriptsize},
        flow/.style={-{Stealth[length=1.7mm]}, color=Logo1, thick}
      ]
        \node[survey-stage] (survey) {Khảo sát trải nghiệm và cảm nhận};
        \node[survey-stage, below=of survey] (indicator) {Tổng hợp chỉ tiêu và lĩnh vực};
        \node[survey-stage, below=of indicator] (province) {Công bố điểm PAPI cấp tỉnh};
        \draw[flow] (survey) -- (indicator);
        \draw[flow] (indicator) -- (province);
      \end{tikzpicture}

      \vspace{0.32cm}
      {\color{Logo1}\Large 18.894}\par
      {\scriptsize người tham gia khảo sát năm 2024}\par
      \vspace{0.22cm}
      {\color{Logo1}\Large 216.673}\par
      {\scriptsize lượt phỏng vấn trực tiếp, 2009-2024}
    \end{column}
  \end{columns}
  \sourcefoot{UNDP Việt Nam, Báo cáo PAPI 2024.}
\end{frame}

% Slide 4
\begin{frame}[t]{Ý nghĩa, đối tượng và mục tiêu phân tích}
  \small
  \begin{columns}[T,onlytextwidth]
    \begin{column}{0.31\textwidth}
      \centerline{\textbf{Ý nghĩa}}
      \begin{itemize}
        \setlength{\itemsep}{7pt}
        \item[-] Theo dõi thay đổi quản trị qua thời gian.
        \item[-] So sánh địa phương trên cùng phạm vi.
        \item[-] Nhận diện lĩnh vực cần ưu tiên xem xét.
      \end{itemize}
    \end{column}
    \begin{column}{0.31\textwidth}
      \centerline{\textbf{Đối tượng sử dụng}}
      \begin{itemize}
        \setlength{\itemsep}{7pt}
        \item[-] Cơ quan quản lý và nhà hoạch định chính sách.
        \item[-] Nhà nghiên cứu, giảng viên và sinh viên.
        \item[-] Người dân quan tâm đến quản trị địa phương.
      \end{itemize}
    \end{column}
    \begin{column}{0.34\textwidth}
      \centerline{\textbf{Mục tiêu phân tích}}
      \begin{itemize}
        \setlength{\itemsep}{7pt}
        \item[-] Diễn biến theo thời gian.
        \item[-] Khác biệt giữa vùng và tỉnh.
        \item[-] Quan hệ giữa các lĩnh vực.
        \item[-] Thay đổi và phân nhóm tỉnh.
      \end{itemize}
    \end{column}
  \end{columns}

  \vspace{0.28cm}
  \centering
  {\footnotesize Dữ liệu PAPI $\rightarrow$ bằng chứng so sánh $\rightarrow$ lĩnh vực ưu tiên $\rightarrow$ định hướng cải thiện phù hợp với địa phương.}
\end{frame}
